\section{模型假设与符号说明}

\subsection{模型假设}

\begin{enumerate}
    \item \textbf{功率无损传输假设}：不计园区内部线路和设备的功率损耗，所有功率传输效率视为100\%，即发电端与用电端功率实时平衡。该假设简化了功率平衡方程，使问题聚焦于能量调配而非损耗计算。\textit{（来源：题目提示"不计园区功率损耗"）}
    
    \item \textbf{稳态运行假设}：园区运行分析时段为1小时，各时段内风电、光伏出力和负荷功率保持恒定，不考虑设备启停时间和爬坡约束。该假设将连续时间问题离散化为24时段整数规划问题，在保证工程精度的前提下显著降低了模型复杂度。\textit{（来源：题目给定分析时段为1小时）}
    
    \item \textbf{设备效率恒定假设}：电解槽制氢效率、合成氨装置产率及储能设备的充放电效率均按额定参数计算，不考虑温度、负载率等因素对效率的影响。该假设使各设备的输入输出关系简化为线性比例关系，便于采用线性规划方法求解。\textit{（来源：题目给定额定参数）}
    
    \item \textbf{成本构成简化假设}：风光发电设备的度电成本已包含其投资折旧和运行维护成本，电解槽和合成氨装置的运行维护成本按耗电量线性计算，不考虑设备折旧对日产氨成本的影响。该假设使吨氨成本的计算聚焦于运行成本，避免了设备投资折旧年限等不确定性因素对结果的影响。\textit{（来源：题目给定度电成本和运维系数）}
    
    \item \textbf{场景等概率假设}：24种风光出力场景等概率出现，每种场景代表15天，全年按360天计。该假设将无限多的实际天气模式归纳为有限个典型场景，在计算可行性与统计代表性之间取得平衡。\textit{（来源：题目给定6种风电场景和4种光伏场景的组合）}
\end{enumerate}

\subsection{符号说明}

本文使用的主要符号及其含义如表\ref{tab:notation}所示，其他变量在文中首次出现时另行说明。

\begin{table}[H]
\caption{主要符号说明}
\label{tab:notation}
\centering
\begin{tabular*}{\textwidth}{@{\extracolsep{\fill}}ccl}
\toprule[1.5pt]
符号 & 单位 & 含义 \\
\midrule[1pt]
$t$ & h & 时段索引，$t=0,1,\ldots,23$ \\
$P_{\mathrm{w}}$ & MW & 风电装机容量（40 MW） \\
$P_{\mathrm{pv}}$ & MW & 光伏装机容量（64 MW） \\
$P_{\mathrm{L}}$ & MW & 常规电负荷（峰值6 MW） \\
$P_{\mathrm{prod}}$ & MW & 制氢氨装置额定功率 \\
$P_{\mathrm{b}}$ & MW & 网购电力功率 \\
$P_{\mathrm{s}}$ & MW & 售电（上网）功率 \\
$C_{\mathrm{b}}$ & 元/MWh & 分时购电价（峰802.4/平607.4/谷342.4） \\
$C_{\mathrm{s}}$ & 元/MWh & 上网电价（377.9） \\
$Q_{\mathrm{NH_3}}$ & t/d & 日制氨产量 \\
$E_{\mathrm{re}}$ & MWh & 新能源发电量 \\
$E_{\mathrm{total}}$ & MWh & 总用电量 \\
$E_{\mathrm{sell}}$ & MWh & 上网电量 \\
$E_{\mathrm{buy}}$ & MWh & 网购电量 \\
$R_{\mathrm{sc}}$ & -- & 新能源自发自用比例（要求$\ge 60\%$） \\
$R_{\mathrm{g}}$ & -- & 总用电绿电比例（要求$\ge 30\%$） \\
$R_{\mathrm{cur}}$ & -- & 新能源上网电量比例（要求$\le 20\%$） \\
\bottomrule[1.5pt]
\end{tabular*}
\end{table}


